% 数论（综述）
% license CCBY4
% type Wiki

本文根据 CC-BY-SA 协议转载翻译自维基百科\href{https://en.wikipedia.org/wiki/Number_theory}{相关文章}

\begin{figure}[ht]
\centering
\includegraphics[width=6cm]{./figures/ae528936f6668ed9.png}
\caption{} \label{fig_SHuLun_1}
\end{figure}
数论是纯数学的一个分支，主要致力于研究整数以及算术函数。数论学家研究素数，以及由整数构造的数学对象（例如有理数）的性质，或作为整数的推广而定义的对象（例如代数整数）。

整数既可以单独研究，也可以作为某些方程的解来研究（丢番图几何）。数论中的问题常常能够通过研究某些解析对象来理解，例如黎曼ζ函数，它以某种方式编码了整数、素数或其他数论对象的性质（解析数论）。另外，还可以研究实数与有理数之间的关系，例如研究无理数如何用分数来逼近（丢番图逼近）。

数论与几何一样，是数学最古老的分支之一。数论的一个特点是，它处理的往往是非常容易理解却极难解决的命题。例如，费马大定理在最初提出后的 358 年才被证明，而哥德巴赫猜想自 18 世纪提出以来至今仍未解决。德国数学家卡尔·弗里德里希·高斯曾经说过：“数学是科学的女王，而数论是数学的女王。”\(^\text{[1]}\)数论一度被视为纯粹数学的典范，认为它在数学之外没有任何应用，直到 20 世纪 70 年代，人们才发现素数可被用作构建公钥密码算法的基础。
\subsection{历史}

数论是数学的一个分支，研究整数及其性质与关系。\(^\text{[2]}\)整数是一个集合，它在自然数集合${1,2,3,\dots}$的基础上扩展，加入了数0，以及自然数的相反数${-1,-2,-3,\dots}$。数论学家不仅研究素数，还研究由整数构造的数学对象（例如有理数）的性质，或定义为整数推广的对象（例如代数整数）。\(^\text{[3][4]}\)

数论与算术密切相关，一些作者甚至将这两个词视为同义词。\(^\text{[5]}\)然而，今天“算术”一词通常指研究数值运算的学科，并扩展到实数范围。\(^\text{[6]}\)在更具体的意义上，数论仅限于研究整数，关注它们的性质和相互关系。\(^\text{[7]}\)传统上，它也被称为高等算术。\(^\text{[8]}\)到 20 世纪初，“数论”这一术语已被广泛采用。\(^\text{[note 1]}\)其中“number”意为整数，既可以指自然数，也可以指全体整数。\(^\text{[9][10][11]}\)

初等数论研究整数的一些方面，这些方面可以用初等方法来探讨，例如利用初等证明。\(^\text{[12]}\)相比之下，解析数论依赖于复数以及来自分析学和微积分的技巧。\(^\text{[13]}\)代数数论则使用代数结构（如域和环）来研究数的性质及其相互关系。几何数论运用几何学的概念研究数。\(^\text{[14]}\)数论的其他分支还包括：概率数论\(^\text{[15]}\)、组合数论\(^\text{[16]}\)、计算数论\(^\text{[17]}\)，以及应用数论，后者研究数论在科学与技术中的应用。\(^\text{[18]}\)
\subsubsection{起源}
\begin{figure}[ht]
\centering
\includegraphics[width=8cm]{./figures/b3c636666bde843e.png}
\caption{普林普顿 322 泥板文书} \label{fig_SHuLun_2}
\end{figure}
在有记录的历史中，数的知识存在于古代的美索不达米亚、埃及、中国和印度文明中\(^\text{[19]}\)。已知最早的算术性质的历史发现是普林普顿 322 号泥板，约公元前 1800 年。它是一块破碎的泥板，包含了一张毕达哥拉斯三元组（即整数$(a,b,c)$，满足$a^{2}+b^{2}=c^{2}$的数对）的列表。这些三元组数量众多且数值很大，不可能通过蛮力枚举得到\(^\text{[20]}\)。

该表的布局表明，它是借助于在现代语言中可写作如下的恒等式构造出来的\(^\text{[21]}\)：
$$
\left(\tfrac{1}{2}\left(x-\tfrac{1}{x}\right)\right)^{2}+1
=\left(\tfrac{1}{2}\left(x+\tfrac{1}{x}\right)\right)^{2},~
$$
这一恒等式隐含在古巴比伦日常的习题中\(^\text{[22]}\)。不过也有人提出，该表可能是作为学校习题的数值例子的来源\(^\text{[23][note 2]}\)。普林普顿 322 泥板是巴比伦数学中唯一现存的证据，可以被视为与现代所谓数论相对应的成果；然而，当时的一种巴比伦代数学要发达得多\(^\text{[24]}\)。

尽管其他文明可能在最初对希腊数学有所影响\(^\text{[25]}\)，但所有关于这种借鉴的证据都出现得相对较晚\(^\text{[26][27]}\)，因此很可能希腊的arithmētikḗ（数的理论性或哲学性研究）是一种本土传统\(^\text{[28]}\)。古希腊人对整除性产生了浓厚兴趣。毕达哥拉斯学派赋予完全数与亲和数神秘的性质。毕达哥拉斯的传统还谈及所谓的多边形数或图形数\(^\text{[29]}\)。欧几里得在其《几何原本》中专门论述了一些属于初等数论的主题，包括素数与整除性\(^\text{[30]}\)。他提出了用于计算两个数的最大公约数的欧几里得算法，并给出了一个证明，表明素数的数量是无限的。晚期古代最重要的arithmētikḗ权威是亚历山大的丢番图，他大约生活在公元三世纪。他撰写了《算术》，这是一部题集，其中的问题无一例外是寻找某个多项式方程组的有理数解，通常形式为：$f(x,y) = z^{2}$或$f(x,y,z) = w^{2}$.用现代术语来说，丢番图方程是指要求有理数或整数解的多项式方程。
\begin{figure}[ht]
\centering
\includegraphics[width=6cm]{./figures/759cf456e7ce5d29.png}
\caption{丢番图《算术》拉丁文译本（1621年，巴谢译）的书名页。} \label{fig_SHuLun_3}
\end{figure}
在罗马帝国灭亡之后，数论的发展转向亚洲，尽管过程断断续续。中国剩余定理出现在《孙子算经》（约三至五世纪）的一道习题中\(^\text{[31][32]}\)。这一结果后来被推广，并在秦九韶1247年的《数学九章》中以完整解法“大衍术”呈现\(^\text{[33][34]}\)。中国数学中也存在某些数秘学思想\(^\text{[note 3]}\)，但与毕达哥拉斯学派不同，这些思想似乎并未带来实质性的数学进展。尽管希腊天文学可能影响了印度的学术\(^\text{[35]}\)，但印度数学在其他方面似乎是一种本土传统\(^\text{[36][37]}\)。阿耶波多（Āryabhaṭa, 476–550 AD）提出，同时同余式对$n \equiv a_{1} \pmod{m_{1}}, \quad n \equiv a_{2} \pmod{m_{2}}$可以通过一种他称为kuṭṭaka（意为“粉碎法”或“逐次约化法”）的算法来求解【38】；这一过程与欧几里得算法十分接近\(^\text{[39]}\)。阿耶波多似乎主要考虑其在天文计算中的应用\(^\text{[35]}\)。婆罗摩笈多（Brahmagupta, 公元628年）则开启了对不定二次方程的系统研究，尤其是Pell方程。一种通解 Pell 方程的方法可能最早由Jayadeva发现；而最早保存下来的系统论述，则见于Bhāskara II（婆什迦罗二世）十二世纪的《种数论》\(^\text{[40]}\)。
\begin{figure}[ht]
\centering
\includegraphics[width=6cm]{./figures/3a3bcc555583aa3f.png}
\caption{西方视野中的伊本·海赛姆：在《月面学》的卷首插图中，‘Alhasen’（即伊本·海赛姆）代表理性之知，而伽利略则代表感官之知。} \label{fig_SHuLun_4}
\end{figure}
九世纪初，哈里发马蒙下令翻译大量希腊数学著作，以及至少一部梵文著作\(^\text{[41][42]}\)。丢番图的主要著作《算术》由古斯塔·伊本·卢卡（Qusta ibn Luqa, 820–912）译为阿拉伯文。卡拉吉（al-Karajī, 953 – 约1029）的部分著作《al-Fakhri》在一定程度上以此为基础。据Rashed Roshdi所言，卡拉吉的同时代人伊本·海赛姆已经知道后来被称为威尔逊定理的结果\(^\text{[43]}\)。在中世纪的西欧，除斐波那契所作关于等差数列平方数的论著外，几乎没有真正意义上的数论研究。直到文艺复兴晚期，随着对希腊古代著作的重新研究，这一局面才开始改变。其催化剂是丢番图《算术》的文本校订和译为拉丁文\(^\text{[44]}\)。
\subsubsection{近代早期数论}
\begin{figure}[ht]
\centering
\includegraphics[width=6cm]{./figures/c7a3992175aa4105.png}
\caption{皮埃尔·德·费马} \label{fig_SHuLun_5}
\end{figure}
法国数学家皮埃尔·德·费马（Pierre de Fermat, 1607–1665）从未正式发表过著作，而是通过通信以及在书页空白处的批注来交流思想\(^\text{[45]}\)。他对数论的贡献重新激发了欧洲对这一领域的兴趣。他提出了费马小定理（模算术中的一个基本结果）、费马大定理，并证明了费马直角三角形定理\(^\text{[2][46]}\)。他还研究了素数、四平方和定理以及Pell 方程\(^\text{[47][48]}\)。莱昂哈德·欧拉（Leonhard Euler, 1707–1783）对数论的兴趣最早可以追溯到 1729 年。当时，他的朋友、业余数学爱好者\(^\text{[note 4]}\)克里斯蒂安·哥德巴赫 向他推荐了费马在这一领域的一些工作\(^\text{[49][50]}\)。在费马未能成功引起同时代学者广泛关注之后，这一事件被称为现代数论的“再生”\(^\text{[51][52]}\)。欧拉证明了费马的一些断言，包括费马小定理；他率先开展了关于“每个整数都可以写成四个平方数之和”的证明工作\(^\text{[53]}\)；并处理了费马大定理的某些特例\(^\text{[54]}\)。他还研究了连分数与 Pell 方程之间的联系\(^\text{[55][56]}\)，并首次迈向了解析数论\(^\text{[57]}\)。
\begin{figure}[ht]
\centering
\includegraphics[width=6cm]{./figures/617824dcffca2fc8.png}
\caption{莱昂哈德·欧拉} \label{fig_SHuLun_6}
\end{figure}
三位欧洲同时代的数学家延续了初等数论的研究：约瑟夫–路易·拉格朗日（Joseph-Louis Lagrange, 1736–1813）给出了四平方和定理、威尔逊定理的完整证明，并发展了 Pell 方程的基本理论。阿德里安–玛丽·勒让德（Adrien-Marie Legendre, 1752–1833）提出了二次互反律，并猜想了相当于素数定理和狄利克雷算术级数定理的结果。他还系统研究了方程
$ ax^{2} + by^{2} + cz^{2} = 0$在他晚年时，他成为第一个证明费马大定理在 (n=5) 情形成立的人\(^\text{[59]}\)。卡尔·弗里德里希·高斯（Carl Friedrich Gauss, 1777–1855）在 1801 年撰写了《算术研究》，该书对数论产生了巨大影响，并为整个 19 世纪的研究奠定了议程。他在书中证明了二次互反律\(^\text{[60]}\)，并发展了二次型理论。他还为同余式引入了一些基本记号，并专门讨论了计算问题，包括素性检验\(^\text{[61]}\)。此外，他建立了单位根与数论之间的联系\(^\text{[62]}\)。通过这些工作，可以说高斯已在某种意义上为后来的埃瓦里斯特·伽罗瓦的研究，以及代数数论领域，开辟了道路。
\subsubsection{数论的成熟与分支划分}
\begin{figure}[ht]
\centering
\includegraphics[width=6cm]{./figures/f56c4c4e4969cc6c.png}
\caption{彼得·古斯塔夫·勒热纳·狄利克雷} \label{fig_SHuLun_7}
\end{figure}
自19世纪早期开始，数论逐渐经历了以下发展：
\begin{itemize}
\item 数论（或高等算术）作为一个独立研究领域的自觉兴起【63】。
\item 现代数学诸多工具的发展，这些工具是基础现代数论所必需的：复分析、群论、伽罗瓦理论——同时伴随着分析的更大 rigor（严格性）以及代数中的抽象化。
\item 数论粗略划分为现代分支，特别是解析数论与代数数论。
\end{itemize}
代数数论可说起始于对互反律与圆分问题的研究，但真正走向成熟则依赖于抽象代数、理想论与赋值论的早期发展（见后文）。解析数论的常规起点是狄利克雷算术级数定理（1837）【64】【65】，其证明引入了L-函数，并涉及一些渐近分析以及对实变量的极限过程【66】。

实际上，最早在数论中使用解析思想可追溯至欧拉（1730年代）【67】【68】，他使用了形式幂级数以及非严格（或隐含）的极限论证。将复分析引入数论则要更晚一些：黎曼（Bernhard Riemann, 1859）关于ζ函数的研究是公认的起点【69】。而 雅可比的四平方和定理（1839）虽更早，却属于最初不同的研究方向，如今已在解析数论中占据核心地位（模形式）【70】。

美国数学会设立了柯尔数论奖。此外，数论还是费马奖所奖励的三大数学分支之一。
\subsection{主要分支}
\subsubsection{初等数论}
\begin{figure}[ht]
\centering
\includegraphics[width=6cm]{./figures/19f35a6cc5cac8ac.png}
\caption{数论学家保罗·埃尔德什与陶哲轩于1985年的合影，当时埃尔德什72岁，陶哲轩10岁} \label{fig_SHuLun_8}
\end{figure}
初等数论通过算术中的基本方法研究数论中的主题。[4] 它的主要研究对象包括整除性、因式分解和素数性，以及模算术中的同余关系。[71][12] 其他的研究内容还包括丢番图方程、连分数、整数拆分以及丢番图逼近。[72]

算术是研究数值运算的学科，探讨数字如何通过加法、减法、乘法、除法、乘方、开方和对数等算术运算进行组合和变换。例如，乘法是一种运算，它将两个数（称为因子）结合成一个数（称为积），如：$2 \times 3 = 6$【73】。

整除性是两个非零整数之间的一种性质，与除法相关。若一个整数$a$可以被一个非零整数 $b$整除，则称$a$是$b$的倍数；换言之，存在某个整数$q$，使得$a = bq$。另一种等价的表述是：$b$整除$a$，记作$b \mid a$。反之，若不满足整除关系，则$a$不能被$b$整除，除法会产生余数。欧几里得的除法引理断言，$a$和$b$一般可以表示为：$a = bq + r$,其中余数$r$是最小的正数剩余量。初等数论研究各种整除规则，以便快速判断一个给定整数是否可以被某个固定的除数整除。例如，众所周知，一个整数如果它的十进制各位数字之和能被 3 整除，那么该整数本身也能被 3 整除。[74][9][75]
\begin{figure}[ht]
\centering
\includegraphics[width=6cm]{./figures/13cc8818feba3d5f.png}
\caption{} \label{fig_SHuLun_9}
\end{figure}
若若干个非零整数有一个公约数，则该整数能够整除它们所有数。最大的这种约数称为最大公约数（gcd）。如果两个整数的最大公约数是 1，即它们唯一的公约数是 1，则称这两个整数互素或相对素数。欧几里得算法可用于计算两个整数 (a, b) 的最大公约数，其方法是反复应用除法引理，并在每一步将除数和余数交换位置。该算法还可以扩展来求解线性丢番图方程的特殊情况：$ax + by = 1$。一个丢番图方程一般具有多个未知数，并且系数都是整数。另一类丢番图方程体现在**勾股定理**中：$x^2 + y^2 = z^2$，如果其解都是整数，则称为勾股三元组。[9][10]另一种重要的表达形式是连分数，它把一个数写成一个整数与一个分数之和，而分母中的分数又是类似形式的整数与分数之和。[76]

初等数论研究整数的整除性质，例如奇偶性（偶数与奇数）、素数和完全数。一些重要的数论函数包括：约数个数函数、约数求和函数及其变形，以及欧拉函数。素数是指大于 1 的整数，且它的正约数只有 1 和它自身。大于 1 的正整数若不是素数，则称为合数。欧几里得定理证明了素数有无穷多个，即集合${2, 3, 5, 7, 11, \dots}$。埃拉托色尼筛法是一种高效算法，它通过逐步剔除所有合数来识别不超过给定自然数的所有素数。[77]

因数分解是一种将一个数表示为乘积的方法。特别是在数论中，整数分解是将一个整数分解为若干整数的乘积。不断重复这一过程直到所有因子都是素数，被称为素因数分解。素数的一个基本性质体现在**欧几里得引理**中。该引理的推论是：如果一个素数能整除若干整数的乘积，那么它必然能整除其中至少一个因子。与素因数分解相关的最基本定理是算术基本定理（又称唯一分解定理）。该定理指出：每个大于 1 的整数都可以分解为素数的乘积，并且这种分解在因子顺序之外是唯一的。例如：$120 = 2 \times 2 \times 2 \times 3 \times 5 = 2^{3} \times 3 \times 5$。[78][9]

模算术处理有限集合中的整数，并引入同余和剩余类的概念。若两个整数$a, b$对模$n$（一个正整数，称为模数）满足$n \mid (a-b)$，则称$a$与$b$在模$n$下同余。即在对$a$与$n$、$b$与$n$进行欧几里得除法时得到相同的余数。这写作：$a \equiv b \pmod{n}$。类似于 12 小时制时钟，$4 + 9 = 13$，但在模 12 的意义下同余于 1。一个模$n$的剩余类是包含所有与某个指定整数$r$在模$n$下同余的整数的集合。例如，$6\mathbb{Z} + 1$包含所有能表示为$6k+1$的整数。

模算术提供了许多公式，可以快速解决涉及大指数幂的同余问题。其中一个重要定理是费马小定理，它指出：如果素数$p$与某个整数$a$互素，则$a^{p-1} \equiv 1 \pmod{p}$。欧拉定理将其推广为：对于任意整数$n$，若$a$与$n$互素，则
$$
a^{\varphi(n)} \equiv 1 \pmod{n}~
$$
其中$\varphi(n)$是欧拉函数，表示小于等于$n$且与$n$互素的正整数的个数。模算术还提供了解决含有未知数的同余方程的公式，类似于代数方程的解法，例如中国剩余定理。[79]
